\chapter{ADMINISTRACIÓN DEL PLAN DE TESIS}
\label{ch:administracion}
\thispagestyle{empty}

Las pruebas numéricas previstas requieren tiempo de cómputo y una secuencia de entregas que permita revisar cada resultado antes de continuar. El cronograma, el presupuesto y el financiamiento se organizan en función de esas necesidades.

\section{CRONOGRAMA}
\label{sec:cronograma}

El estudio se programa en doce meses, de julio de 2026 a junio de 2027. El corte de esta revisión es el 17 de noviembre de 2026, posterior al informe del asesor del 13 de noviembre. La Tabla~\ref{tab:cronograma} distingue actividad documentada y programación. El bloque julio--septiembre conserva la planificación inicial, sin acreditar fechas de cierre que no tengan acta o evidencia fechada.

Las campañas de octubre corresponden al desarrollo exploratorio. La confirmación del protocolo revisado se programa para noviembre--diciembre y depende de superar el piloto.

\providecommand{\gc}{\cellcolor{darkblue}}
\providecommand{\gp}{\cellcolor{darkblue!35}}
\begin{table}[htbp]
\centering
\caption{Cronograma revisado al 17 de noviembre de 2026. Azul oscuro: actividad documentada, no necesariamente aceptada; azul claro: programación inicial o pendiente. Fuente: Elaboración propia.}
\label{tab:cronograma}
\scriptsize
\setlength{\tabcolsep}{2.5pt}
\begin{tabular}{C{0.7cm}L{4.3cm}*{10}{C{0.82cm}}}
\toprule
& & & \multicolumn{3}{c}{\textbf{2026}} & \multicolumn{6}{c}{\textbf{2027}} \\
\cmidrule(lr){4-6}\cmidrule(lr){7-12}
\textbf{Fase} & \textbf{Actividad} & \textbf{\makecell{Jul--\\sep}} & \textbf{Oct} & \textbf{Nov} & \textbf{Dic} & \textbf{Ene} & \textbf{Feb} & \textbf{Mar} & \textbf{Abr} & \textbf{May} & \textbf{Jun} \\
\midrule
1 & Revisión y cierre de especificación & \gp &  & \gc & \gp &  &  &  &  &  &  \\
2 & Diseño y prototipo; ajuste al protocolo & \gp & \gc & \gc & \gp &  &  &  &  &  &  \\
3 & Desarrollo exploratorio y piloto & & \gc & \gc & \gp & & & & & & \\
3 & Confirmación F2 con semillas nuevas &  &  & \gp & \gp &  &  &  &  &  &  \\
4 & Demostración y evaluación térmica F3 & {} & {} & {} & \gp & \gp & {} & {} & {} & {} & {} \\
5 & Comparación de ampacidad F4 &  &  &  &  & \gp & \gp &  &  &  &  \\
6 & Síntesis DSR &  &  &  &  &  & \gp & \gp &  &  &  \\
7 & Comunicación: redacción, revisión del asesor y sustentación &  &  &  &  &  &  & \gp & \gp & \gp & \gp \\
\bottomrule
\end{tabular}
\end{table}

La bitácora del proyecto registra campañas adaptativas del 26--27 de octubre de 2026 con 39, 38, 40 y 42 corridas, y una extensión del 27--28 de octubre con 22. Su identificación y clasificación se conservan en los informes de campañas del repositorio \texttt{pinnV3} del 27 de octubre de 2026. Son antecedentes de desarrollo: las reglas históricas de parada, aceptación y uso de referencias difieren del protocolo revisado. No acreditan H1 ni H2, ni permiten declarar una arquitectura ganadora para la confirmación.

Las pruebas de noviembre inventariadas en la bitácora son exploratorias, parciales o de funcionamiento; no constituyen confirmación independiente. Al corte de revisión no se declara cumplido ningún hito de aceptación con esas ejecuciones. La ronda E deja de ser un compromiso separado: cualquier desarrollo adicional dependerá del fallo documentado del piloto y se versionará antes de la confirmación. La Tabla~\ref{tab:hitos-revisados} fija la evidencia exigida para cerrar cada hito.

\begin{table}[htbp]
\centering
\caption{Hitos reprogramados y evidencia exigida para su cierre. Fuente: Elaboración propia.}
\label{tab:hitos-revisados}
\begin{tabularx}{\textwidth}{L{1.1cm}L{2.4cm}Y}
\toprule
\textbf{Hito} & \makecell[l]{\textbf{Fecha}\\\textbf{prevista}} & \textbf{Estado al corte y evidencia de cierre} \\
\midrule
H1 & Nov.--dic. 2026 & En revisión. Para su cierre se requieren VD1, fichas, matriz requisito--prueba, entorno versionado y estudio de convergencia FEM aprobados. \\
H2 & Dic. 2026 & Pendiente. Manifiesto congelado, cinco semillas nuevas por caso F2 y expediente VD2 con errores, controles físicos y fallos. \\
H3 & Nov.--dic. 2026 & Pendiente. Acta de decisión del piloto: confirmar o rediseñar; sin asignar aceptación a la antigua ronda E. \\
H4 & Ene.--feb. 2027 & Pendiente. VD3 y VD4, pares completos, intervalos de ampacidad y clasificación según incertidumbre. \\
H5 & Mar. 2027 & Pendiente. Borrador íntegro, repositorio reproducible y síntesis DSR. \\
H6 & Abr.--jun. 2027 & Pendiente. Revisión del asesor, versión de depósito y sustentación conforme a requisitos institucionales. \\
\bottomrule
\end{tabularx}
\end{table}

Si H1--H3 no se cierran en diciembre, se reprogramarán las fases dependientes con el asesor. No se sustituirán casos difíciles por resultados favorables ni se ampliarán umbrales para conservar una fecha de entrega.

\section{PRESUPUESTO}
\label{sec:presupuesto}

La Tabla~\ref{tab:presupuesto} reúne los gastos previstos en soles para los doce meses de estudio. El equipo ya está disponible y el software es de código abierto; no se prevén compras por esos conceptos. Los demás importes son estimaciones pendientes de contrastar con cotizaciones y tarifas oficiales.

Para conectividad y energía se asignan montos mensuales; los derechos de sustentación y grado se cubren con una provisión. Estos valores se comprobarán antes de asumir cada gasto y los comprobantes se conservarán durante la ejecución.

\begin{table}[htbp]
\centering
\caption{Presupuesto estimado de planificación, en soles. Fuente: Elaboración propia; importes sujetos a contraste antes del gasto.}
\label{tab:presupuesto}
\scriptsize
\begin{tabularx}{\textwidth}{L{3.0cm}Y L{2.5cm}L{2.2cm}}
\toprule
\textbf{Rubro} & \textbf{Concepto} & \textbf{Base estimada} & \textbf{Monto (S/)} \\
\midrule
a) Equipos & Computadora personal con CPU Intel Core Ultra 7 165U para entrenamientos y referencias FEM. & Existente, de los tesistas & 0,00 \\
b) Materiales e insumos & Software de código abierto: FEniCSx, Gmsh, PETSc, PyTorch y LaTeX. & Sin costo de licencia & 0,00 \\
b) Materiales e insumos & Impresión y encuadernación. & 3 ejemplares a S/ 80 & 240,00 \\
c) Servicios & Asignación de conectividad. & 12 meses a S/ 30 & 360,00 \\
c) Servicios & Asignación de energía eléctrica. & 12 meses a S/ 25 & 300,00 \\
c) Servicios & Provisión para derechos de sustentación y grado; no representa el tarifario UNI. & Reserva estimada & 1500,00 \\
c) Servicios & Reserva de capacidad de cómputo externa si el piloto lo requiere. & Tope de planificación & 300,00 \\
d) Viáticos y publicación & Movilidad para asesoría y sustentación. & 12 traslados a S/ 20 & 240,00 \\
d) Viáticos y publicación & Preparación y envío del manuscrito a una revista sin cargo de publicación. & Ruta sin APC & 0,00 \\
\midrule
\multicolumn{3}{r}{Subtotal estimado} & 2940,00 \\
\multicolumn{3}{r}{Contingencia: 20\,\pct{} del subtotal} & 588,00 \\
\multicolumn{3}{r}{\textbf{Total estimado}} & \textbf{3528,00} \\
\bottomrule
\end{tabularx}
\end{table}

La publicación depende de la evaluación editorial y no se garantiza aceptación. Si se elige una revista con APC o el tarifario institucional supera la provisión, se actualizarán presupuesto y financiamiento antes de asumir el compromiso. El costo del equipo existente no se trata como desembolso nuevo. La estimación financia el estudio, no es una evaluación económica del cable o de la instalación.

\subsection{RIESGOS Y CONTINGENCIAS}

La Tabla~\ref{tab:riesgos-plan} vincula los riesgos de ejecución con acciones y reglas de continuidad. La revisión mensual permitirá actualizar plazos y recursos sin modificar retrospectivamente los criterios de aceptación.

\begin{table}[htbp]
\centering
\caption{Riesgos principales y respuesta prevista. Fuente: Elaboración propia.}
\label{tab:riesgos-plan}
\begin{tabularx}{\textwidth}{L{3.2cm}Y Y}
\toprule
\textbf{Riesgo / señal} & \textbf{Prevención y respuesta} & \textbf{Regla de continuidad} \\
\midrule
Cómputo o memoria insuficientes & Medir tiempo y memoria del piloto, guardar estados y usar CPU/WSL si el acelerador no admite precisión o determinismo requeridos. & Extrapolar el costo al catálogo y a la bisección antes de confirmar; usar la reserva solo si basta, o reprogramar con el asesor. \\
No convergencia PINN o electrotermal & Revisar escalamiento, interfaces y pérdidas en desarrollo; conservar errores y semillas fallidas. & No relajar puertas; rediseñar, versionar y repetir la confirmación completa del dominio afectado. \\
Referencia FEM no convergente & Refinar malla, cuadratura y dominio; verificar primero casos exactos. & No aceptar el caso sin referencia estable; documentar la restricción del dominio. \\
Datos publicados incompletos & Auditar cada dato y distinguir reproducción de adaptación o caso sintético. & Mantener el catálogo mínimo explícito; no completar datos faltantes atribuyéndolos al artículo. \\
Pérdida de archivos o interrupción & Versionar código y manifiestos, respaldar resultados y registrar huellas y reanudaciones. & Repetir ejecuciones no reconstruibles antes de usarlas como evidencia. \\
Retrasos o costos superiores & Revisar avance mensual, tarifas y disponibilidad de expertos antes de su fase. & Reprogramar y actualizar la provisión; el juicio experto no reemplazará la verificación numérica. \\
\bottomrule
\end{tabularx}
\end{table}

\section{FINANCIAMIENTO}
\label{sec:financiamiento}

Los tesistas prevén cubrir con fondos propios el total estimado de S/ 3528,00, incluida la contingencia. No se han solicitado fondos concursables ni financiamiento de terceros, y no se requiere un compromiso económico adicional de la Universidad. La continuidad del estudio dependerá del cierre del piloto, la evidencia de cada hito y la actualización de los costos antes de ejecutar los gastos.

